\section{RESULTADOS E DISCUSSÕES}
Esta seção apresenta os procedimentos de verificação e os resultados obtidos no desenvolvimento do Alchemon. São descritos os ambientes de teste, os critérios de desempenho, a verificação dos requisitos funcionais e as principais limitações identificadas.


\subsection{Ambiente de Testes Funcionais e Verificação de Requisitos}
O desenvolvimento foi realizado em ambiente Linux, na \texit{engine} Godot 4.4.1. [Para futuras verificações, estão previstos ambientes Windows e Linux, cujas configurações deverão ser registradas conforme a Tabela I. Os testes de desempenho ainda não foram realizados; portanto, os valores apresentados na Tabela II representam critérios definidos para avaliações futuras, e não resultados obtidos.]
\begin{table}[!t]
    \caption{Ambientes previstos para os testes}
    \label{tab:ambientes-teste}
    \centering
    \footnotesize
    \begin{tabular}{|l|c|c|}
        \hline
        \textbf{Configuração} & \textbf{Windows} & \textbf{Linux} \\ \hline
            Sistema operacional & A preencher & A preencher \\ \hline
            Versão do sistema & A preencher & A preencher \\ \hline
            Versão do Godot & A preencher & A preencher \\ \hline
            Processador & A preencher & A preencher \\ \hline
            Memória RAM & A preencher & A preencher \\ \hline
            Placa de vídeo & A preencher & A preencher \\ \hline
            Resolução & A preencher & A preencher \\ \hline
            Data do teste & A preencher & A preencher \\ \hline
            Situação & Pendente & Pendente \\ \hline
    \end{tabular}
\end{table}

\begin{table}[!t]
    \caption{Critérios definidos para avaliação de desempenho}
    \label{tab:criterios-desempenho}
    \centering
    \footnotesize
    \begin{tabular}{|p{0.23\columnwidth}|p{0.43\columnwidth}|p{0.20\columnwidth}|}
        \hline
        \textbf{Métrica} & \textbf{Critério} & \textbf{Resultado} \\ \hline
            Taxa de quadros & Média igual ou superior a 60 FPS, sem quedas sustentadas abaixo de 30 FPS. & Não mensurado \\ \hline
            Carregamento & Média de até 3 segundos, sem ocorrências acima de 5 segundos. & Não mensurado \\ \hline
            Memória RAM e VRAM & Ausência de crescimento contínuo após o carregamento e cinco ciclos de combate. & Não mensurado \\ \hline
    \end{tabular}
\end{table}

\subsection{Matriz de Verificação de Requisitos Funcionais}
Foram desenvolvidas quatro suítes de testes automatizados com o GdUnit4, totalizando 65 casos. A matriz da Tabela III relaciona os principais requisitos contemplados. A implementação de um teste não implica sua execução ou aprovação; esses estados devem ser registrados após a realização das verificações correspondentes.

\begin{table}[!t]
    \caption{Matriz de verificação dos requisitos funcionais}
    \label{tab:matriz-requisitos}
    \centering
    \footnotesize
    \begin{tabular}{|p{0.27\columnwidth}|p{0.25\columnwidth}|p{0.18\columnwidth}|p{0.18\columnwidth}|}
        \hline
        \textbf{Requisito} & \textbf{Teste implementado} & \textbf{Executado} & \textbf{Integração confirmada} \\ \hline
            Crescimento das criaturas & Sim & A registrar & A registrar \\ \hline
            Atributos das criaturas & Sim & A registrar & A registrar \\ \hline
            Captura & Sim & A registrar & A registrar \\ \hline
            Distribuição de recompensas & Sim & A registrar & A registrar \\ \hline
            Interface & Sim & A registrar & A registrar \\ \hline
            Persistência de dados & Sim & A registrar & A registrar \\ \hline
        \end{tabular}
\end{table}
\subsubsection{Resultados Alcançados}
Durante o desenvolvimento, foram implementados os sistemas de exploração, combate por turnos, captura, progressão e persistência de dados. Também foram identificados e corrigidos problemas relacionados à sincronização dos dados das criaturas durante o combate, à distinção entre atributos temporários e permanentes e ao encerramento das batalhas com a distribuição de recompensas.

Apesar desses avanços, a execução e a aprovação de todos os testes ainda precisam ser confirmadas. Também não foram realizadas avaliações com estudantes, de modo que os efeitos do jogo sobre o engajamento e a aprendizagem ainda não podem ser afirmados.

\subsection{Discussão Crítica e Limitações}

Os resultados obtidos demonstram avanços na implementação das estruturas técnicas do Alchemon, especialmente nos sistemas de exploração, combate, captura, progressão e persistência. Entretanto, o protótipo ainda não apresenta todas as condições necessárias para uma avaliação completa de seu funcionamento e de sua proposta educacional.

Uma das limitações corresponde à ausência de medições de desempenho. Embora tenham sido definidos critérios para avaliar a taxa de quadros por segundo, o tempo de carregamento das cenas e o consumo de memória, esses parâmetros ainda não foram mensurados. Além disso, a execução nativa dos artefatos exportados para Windows x86-64 não foi confirmada. Dessa forma, a presença dos arquivos de exportação não permite concluir que o jogo funciona integralmente nessa plataforma.

Outra limitação está relacionada à abrangência dos testes funcionais. O projeto possui 65 casos de teste automatizados distribuídos em quatro suítes, mas a quantidade de testes implementados não comprova que todos foram executados e aprovados. Também é necessário distinguir a verificação isolada das regras de negócio da confirmação de seu funcionamento integrado à interface e aos demais sistemas do jogo.

Por fim, não foram realizadas avaliações com estudantes do Ensino Médio. Assim, ainda não existem dados empíricos que permitam determinar se as mecânicas de captura e colecionismo favorecem o engajamento ou a aprendizagem tangencial dos conteúdos relacionados à Tabela Periódica. Aspectos como compreensão das mecânicas, curva de aprendizado, nível de dificuldade e experiência de uso permanecem sem validação junto ao público-alvo.

Diante dessas limitações, a continuidade do projeto exige a conclusão e integração das funcionalidades planejadas, a execução dos testes nos ambientes Windows e Linux e a realização das medições de desempenho. Posteriormente, será possível conduzir avaliações com estudantes, incluindo testes de usabilidade e instrumentos de avaliação da aprendizagem. Essas etapas permitirão investigar tanto a estabilidade técnica do jogo quanto sua contribuição para os objetivos educacionais propostos.

\section{CONCLUSÕES}
Neste trabalho, foi apresentado o desenvolvimento do Alchemon, um jogo digital educativo concebido com o propósito de promover a aprendizagem tangencial e o engajamento de estudantes do Ensino Médio com os conceitos e elementos da Tabela Periódica. Ao longo do projeto, foram implementadas estruturas fundamentais para o funcionamento do protótipo, incluindo sistemas de exploração, combate, captura de criaturas, progressão e persistência de dados. Os testes internos, organizados em quatro conjuntos com 65 casos de teste, permitiram verificar diferentes comportamentos dessas estruturas, identificar falhas e orientar suas correções. Dessa forma, os resultados demonstram o cumprimento de objetivos técnicos importantes estabelecidos para esta etapa do trabalho, embora o desenvolvimento da solução ainda não esteja concluído.

Como principal limitação, destaca-se que a validação realizada se restringiu aos aspectos técnicos e funcionais contemplados pelos testes internos. Não foram realizados testes de usabilidade ou avaliações com estudantes, o que impossibilita determinar, neste momento, se o jogo favorece efetivamente a aprendizagem dos conteúdos de Química. Aspectos como a curva de aprendizado, a compreensão das mecânicas, o nível de dificuldade, o engajamento e a recepção dos jogadores ainda dependem de investigação empírica. Portanto, é necessário distinguir o funcionamento das funcionalidades verificadas da eficácia pedagógica da proposta, que permanece sem comprovação prática.

Como contribuição, o trabalho apresenta uma implementação inicial de um jogo educativo que utiliza mecânicas de captura, coleção e desenvolvimento de criaturas como recursos para abordar os elementos da Tabela Periódica. O protótipo e sua arquitetura constituem uma base técnica para a continuidade do projeto e para futuras investigações sobre a aplicação de mecânicas de jogos de captura e colecionismo no ensino de Química. Embora ainda não esteja em condições adequadas para aplicação com o público-alvo, os resultados obtidos demonstram a viabilidade de dar continuidade à solução sem a necessidade de reduzir seu escopo original.

\subsection{Trabalhos Futuros}

Como trabalhos futuros, o projeto demanda a continuidade do desenvolvimento, com tempo adicional para concluir as funcionalidades planejadas e integrar os diferentes sistemas do jogo. Após essa etapa, será necessário realizar uma nova rodada de validação técnica, contemplando a execução dos testes internos, a verificação dos builds para os sistemas operacionais previstos e a avaliação da estabilidade do jogo. Com uma versão suficientemente completa e estável, poderão ser conduzidas avaliações com estudantes do Ensino Médio, por meio de sessões de playtest destinadas a investigar aspectos como usabilidade, dificuldade, engajamento e aprendizagem dos conteúdos relacionados à Tabela Periódica. Essas avaliações também permitirão analisar se as mecânicas de captura, coleção e interação com as criaturas contribuem para a aprendizagem tangencial dos conceitos químicos. Portanto, a continuidade do projeto depende de mais tempo de desenvolvimento e validação, antes de sua avaliação junto ao público-alvo.
